\documentclass[11pt]{article}
\usepackage[a4paper,margin=25mm]{geometry}
\usepackage{amsmath,amssymb,amsthm,booktabs,microtype}
\usepackage[T1]{fontenc}
\usepackage{lmodern}
\usepackage[dvipsnames]{xcolor}
\usepackage[colorlinks=true,linkcolor=MidnightBlue,citecolor=MidnightBlue,urlcolor=MidnightBlue]{hyperref}
\usepackage{fancyhdr}
\pagestyle{fancy}\fancyhf{}
\fancyhead[L]{\small Three-uniform Berge-path saturation}
\fancyhead[R]{\small Dev D. Goyal}
\fancyfoot[C]{\thepage}
\setlength{\headheight}{14pt}
\hypersetup{pdftitle={Exact three-uniform Berge-path saturation numbers and a counterexample to a conjectured formula},
pdfauthor={Dev D. Goyal},
pdfsubject={Fully autonomous agentic research},
pdfkeywords={Berge hypergraphs, Berge paths, saturation numbers, three-uniform hypergraphs, loose cycles, Lean 4}}
\newtheorem{theorem}{Theorem}
\newtheorem{lemma}[theorem]{Lemma}
\newtheorem{proposition}[theorem]{Proposition}
\newtheorem{corollary}[theorem]{Corollary}
\theoremstyle{definition}
\newtheorem{definition}[theorem]{Definition}
\theoremstyle{remark}
\newtheorem{remark}[theorem]{Remark}
\newcommand{\sat}{\operatorname{sat}}
\newcommand{\BP}{\text{Berge-}P}
\title{\textbf{Exact three-uniform Berge-path saturation numbers\\and a counterexample to a conjectured formula}}
\author{Dev D. Goyal}
\date{31 October 2026}
\begin{document}
\maketitle
\begin{abstract}
For $m\ge 10$ write $m=4s+r$ with $1\le r\le 4$, and put $d=m-2$,
$a=(3+r)2^s-5$ and $b=2a+1$. English, Graber, Kirkpatrick, Methuku and
Sullivan conjectured that the three-uniform saturation number of the
Berge path $P_m$ equals $U(n)=\lceil (n-\lfloor n/b\rfloor)/2\rceil$ for every
$n\ge b$. Two disjoint loose $8$-cycles together with two isolated
vertices form a Berge-$P_{10}$-saturated three-uniform hypergraph on $34$
vertices with $16$ hyperedges, one fewer than the conjectured value $17$.
We determine the saturation number exactly on the whole conjectured range:
for all $m\ge10$ and $n\ge b$, $\sat_3(n,\BP_m)=U(n)-1$ when $b<n<2b$ and
$n\equiv 2\pmod{2d}$, and $\sat_3(n,\BP_m)=U(n)$ otherwise. At every
exceptional order the extremal hypergraph is unique: disjoint loose
$d$-cycles and two isolated vertices. The lower bound uses two results of
the source paper, the minimum size of a saturated linear tree and the
general upper construction. Lean~4 verifies, for actual hypergraphs, Berge
paths and unrestricted saturation, that disjoint loose $d$-cycles with two
isolated vertices are Berge-$P_{d+2}$-saturated for every $d\ge3$, that the
explicit $34$-vertex hypergraph is Berge-$P_{10}$-saturated with $16$
hyperedges, and that the three-uniform case of the conjecture fails for
every $m\ge10$. The matching lower bound and the uniqueness statement have
conventional proofs.
\end{abstract}

\section{The conjecture and the result}

All hypergraphs are finite and three-uniform; the vertex set includes
isolated vertices. A \emph{Berge path of length $L$} in a hypergraph $H$
consists of distinct vertices $v_0,\dots,v_L$ and distinct hyperedges
$e_1,\dots,e_L$ of $H$ with $\{v_{i-1},v_i\}\subseteq e_i$ for every $i$.
The graph path $P_m$ has $m$ vertices, so a Berge-$P_m$ is a Berge path of
length $m-1$. A hypergraph $H$ is \emph{Berge-$P_m$-saturated} if it
contains no Berge-$P_m$ but $H+f$ contains one for every three-element set
$f\notin E(H)$; no condition such as linearity is imposed on $f$. The
saturation number $\sat_3(n,\BP_m)$ is the least number of hyperedges of a
Berge-$P_m$-saturated hypergraph on $n$ vertices~\cite{EGKMS}.

A \emph{linear tree} is a connected hypergraph whose incidence graph is a
tree; it has $2e+1$ vertices when it has $e$ hyperedges. Fix $m\ge 10$,
write $m=4s+r$ with $1\le r\le4$ and put
\[
d=m-2,\qquad a=(3+r)2^s-5,\qquad b=2a+1,\qquad
U(n)=\left\lceil\frac{n-\lfloor n/b\rfloor}{2}\right\rceil .
\]
We use two results of~\cite{EGKMS} for $k=3$.
\begin{itemize}
\item[(T)] (Theorem~2.2.) Every Berge-$P_m$-saturated linear tree on at
least four vertices has at least $a$ hyperedges, and one with exactly $a$
hyperedges exists.
\item[(C)] (Construction~3.8, Observation~3.9 and Lemma~3.10.) For every
$n\ge b$ there is a Berge-$P_m$-saturated hypergraph on $n$ vertices with
$U(n)$ hyperedges.
\end{itemize}
Theorem~2.3 of~\cite{EGKMS} gives a lower bound that differs from $U(n)$ by
at most one, and Conjecture~2.4 of~\cite{EGKMS} asserts, for $k=3$, that
$\sat_3(n,\BP_m)=U(n)$ for all $m\ge10$ and $n\ge b$. The authors note that
the bounds would agree if a minimal saturated hypergraph could be shown to
have no isolated vertices. Isolated vertices are exactly what breaks the
formula.

\begin{theorem}\label{thm:main}
Let $m\ge10$ and $n\ge b$. Then
\[
\sat_3(n,\BP_m)=
\begin{cases}
U(n)-1, & \text{if } b<n<2b \text{ and } n\equiv 2 \pmod{2d},\\
U(n), & \text{otherwise.}
\end{cases}
\]
At an exceptional order $n=2td+2$ every extremal hypergraph is isomorphic
to $t$ disjoint loose $d$-cycles together with two isolated vertices.
\end{theorem}

\begin{corollary}\label{cor:counter}
The three-uniform case of Conjecture~2.4 of~\cite{EGKMS} is false for every
$m\ge10$. For $m=10$ and $n=34$ the conjectured value is $17$, while two
disjoint loose $8$-cycles with two isolated vertices give a
Berge-$P_{10}$-saturated hypergraph with $16$ hyperedges.
\end{corollary}

The novel parts are the exceptional construction, the treatment of
isolated vertices and the resulting matching lower bound; results~(T)
and~(C) are inherited. Corollary~\ref{cor:counter} needs only
Proposition~\ref{prop:cycles} below and arithmetic, and it is verified in
Lean (Section~\ref{sec:lean}).

\section{Loose cycles with two isolated vertices}

For $d\ge3$, a \emph{loose $d$-cycle} has core vertices $c_0,\dots,c_{d-1}$,
pendant vertices $p_0,\dots,p_{d-1}$ and hyperedges
$e_i=\{c_i,c_{i+1},p_i\}$, indices modulo~$d$. Let $\mathcal C_{d,t}$
consist of $t$ disjoint loose $d$-cycles and two isolated vertices
$z_1,z_2$; it has $2td+2$ vertices and $td$ hyperedges.

\begin{proposition}\label{prop:cycles}
For $d\ge3$ and $t\ge0$, $\mathcal C_{d,t}$ is Berge-$P_{d+2}$-saturated.
\end{proposition}

\begin{proof}
A Berge path is connected through its hyperedges, so all its hyperedges
lie in one component; each component has at most $d$ hyperedges, while a
Berge-$P_{d+2}$ needs $d+1$. Hence $\mathcal C_{d,t}$ is Berge-$P_{d+2}$-free.

Every vertex of a cycle starts a Berge path through all $d$ hyperedges of
its cycle: from $c_i$ use $c_i,c_{i+1},\dots,c_{i+d-1},p_{i+d-1}$ with
$e_i,\dots,e_{i+d-1}$, and from $p_i$ use the reverse of the path
$c_{i+1},\dots,c_{i+d},p_{i+d}=p_i$. Let $f$ be an absent triple.

If $f$ contains a vertex $u$ of some cycle and a vertex $w$ outside that
cycle, prepend $w,f$ to the $d$-edge path from $u$. Since there are only
two isolated vertices, this covers every $f$ that is not contained in a
single cycle.

If $f$ lies in one cycle, two of its vertices have the same type. If
$p_i,p_j\in f$ with $i\ne j$, prepend $p_j,f$ to the $d$-edge path from
$p_i$, which contains no other pendant vertex. If $c_x,c_y\in f$, let
$k\in\{1,\dots,d-1\}$ with $y\equiv x+k$. Take
\[
p_{x+k-1},\,c_{x+k-1},\,c_{x+k-2},\dots,c_x,\;\;c_y,\,c_{y+1},\dots,c_{y+d-k-1},\,p_{y+d-k-1}
\]
with the hyperedges $e_{x+k-1},\dots,e_x$, then $f$, then
$e_y,\dots,e_{y+d-k-1}$. The two runs of cores occupy the disjoint
offset ranges $[0,k-1]$ and $[k,d-1]$ from $x$, the pendants $p_{x+k-1}$
and $p_{x-1}$ are distinct, and every old hyperedge occurs once. This is a
Berge path with $k+1+(d-k)=d+1$ hyperedges.
\end{proof}

\section{Incidence counting and isolated hyperedges}

For a component with $v$ vertices and $e$ hyperedges, the incidence graph
has $v+e$ nodes and $3e$ edges, hence cyclomatic number $2e-v+1$. A linear
tree has $v=2e+1$; every other component has $v\le2e$, with equality
exactly for incidence-unicyclic components. Berge paths correspond to
paths of the incidence graph between vertex nodes.

Let $H$ be Berge-$P_m$-saturated on $n\ge b$ vertices. It has at most two
isolated vertices, since a triple of isolated vertices could be added
without creating a path of length $d+1\ge9$.

\begin{lemma}\label{lem:edge}
If $H$ has a component consisting of one hyperedge, then
$|E(H)|\ge\lceil (n-1)/2\rceil\ge U(n)$.
\end{lemma}

\begin{proof}
Let $\{x,y,w\}$ be that component and suppose that $T$ is another linear
tree component, possibly a single vertex. Since $H$ is
Berge-$P_m$-free, Berge paths in $T$ have length at most $d$, so its
incidence tree has diameter at most $2d$. A centre of this tree, replaced by
an adjacent vertex node if it is a hyperedge node, gives a vertex $v$ of
$T$ from which every Berge path in $T$ has length at most $\lceil d/2\rceil$.
A Berge path in $H+\{x,y,v\}$ that uses the new hyperedge uses at most one
hyperedge of $\{x,y,w\}$ and at most $\lceil d/2\rceil$ of $T$, so it has
length at most $\lceil d/2\rceil+2\le d$, contradicting saturation. Hence every
other component has $v\le2e$, and $|E(H)|\ge 1+(n-3)/2$. Since
$\lfloor n/b\rfloor\ge1$, the bound is at least $U(n)$.
\end{proof}

\section{One isolated vertex}

\begin{lemma}\label{lem:height}
Let $z$ be an isolated vertex of $H$, let $T$ be a linear tree component
with at least one hyperedge and let $e=\{x,y,w\}$ be a hyperedge of $T$.
Removing $e$ splits $T$ into three branches containing $x$, $y$ and $w$;
let $h_x,h_y,h_w$ be the largest lengths of Berge paths in these branches
starting at $x$, $y$, $w$, labelled so that $h_x\ge h_y\ge h_w$. Then
$h_y=h_w$ and $h_x+h_y=d-1$.
\end{lemma}

\begin{proof}
Joining longest paths in the $x$- and $y$-branches through $e$ gives a
path of length $h_x+h_y+1$, so $h_x+h_y\le d-1$. Add $f=\{z,x,y\}$. A new
path uses $f$. If it avoids $e$, it either joins the $x$- and $y$-branches
through $f$ or ends at $z$, and has length at most $h_x+h_y+1\le d$. If it
uses $e$ and $f$, these must be consecutive with a common vertex
$c\in\{x,y\}$ between them, because the branches are joined only through
$e$ and $f$. The remaining ends are then one of $w$ or the other element of
$\{x,y\}$ on the side of $e$, and one of $z$ or the other element of
$\{x,y\}$ on the side of $f$; checking the four possibilities gives length
at most $h_x+h_w+2$. Saturation forces $h_x+h_w\ge d-1$, and with
$h_w\le h_y$ both equalities follow.
\end{proof}

In particular every such $T$ has Berge diameter exactly $d$.

\begin{lemma}\label{lem:tree}
Under the hypotheses of Lemma~\ref{lem:height}, $T$ has at least $B_m$
hyperedges, where $B_m=2(2^h-1)$ if $d=2h$ and $B_m=3\cdot2^h-2$ if
$d=2h+1$. Moreover $B_m\ge 2a$ for all $m\ge10$.
\end{lemma}

\begin{proof}
The incidence tree has diameter $2d$, and its centre is a vertex node when
$d$ is even and a hyperedge node when $d$ is odd.

Let $d=2h$ with centre $c$. Root the tree at $c$; at least two branches at
$c$ reach Berge depth $h$, and no vertex has depth more than $h$. For a
hyperedge whose parent vertex has depth $j$, the parent side has height
exactly $h+j$ (through another deep branch at $c$), and both child sides
have height at most $h-j-1$. Lemma~\ref{lem:height} forces both child
heights to equal $h-j-1$. Hence every child vertex of depth less than $h$
has a child hyperedge, and a deep branch at $c$ contains at least
$1+2+\dots+2^{h-1}=2^h-1$ hyperedges. Two deep branches give $2(2^h-1)$.

Let $d=2h+1$ with central hyperedge $e_0$. Its three branches have height
at most $h$, and Lemma~\ref{lem:height} applied to $e_0$ makes all three
equal to $h$. Below a vertex of $e_0$, a hyperedge whose parent is $j$
steps from that vertex has parent side of height $h+j+1$ and child sides of
height at most $h-j-1$, which Lemma~\ref{lem:height} again makes equal.
Each branch therefore has at least $2^h-1$ hyperedges, and $T$ has at least
$1+3(2^h-1)$.

For $m=10,11,12,13$ the values $(a,B_m)$ are $(15,30)$, $(19,46)$, $(23,62)$
and $(27,94)$. Moreover $a_{m+4}=2a_m+5$ and $B_{m+4}=4B_m+6$, so
$B_m\ge2a_m$ implies $B_{m+4}-2a_{m+4}=4(B_m-a_m-1)\ge0$.
\end{proof}

\begin{corollary}\label{cor:one}
If $H$ has exactly one isolated vertex, then $|E(H)|\ge U(n)$.
\end{corollary}

\begin{proof}
By Lemma~\ref{lem:edge} we may assume that no component is a single
hyperedge. Let $t$ be the number of linear tree components with a
hyperedge; by Lemma~\ref{lem:tree} each has at least $2B_m+1\ge 2b-1$
vertices, so $n\ge 1+t(2b-1)\ge(t+1)b$ when $t\ge1$, and
$\lfloor n/b\rfloor\ge t+1$ in all cases. Counting vertices over the
non-isolated components gives $2|E(H)|\ge n-1-t\ge n-\lfloor n/b\rfloor$.
\end{proof}

\section{Two isolated vertices}

\begin{lemma}\label{lem:long}
If $H$ has two isolated vertices $z_1,z_2$, then every non-isolated vertex
$v$ starts a Berge path of length $d$, and no vertex node is a cut node of
the incidence graph of its component.
\end{lemma}

\begin{proof}
A path created by adding $\{z_1,z_2,v\}$ must end with this hyperedge at
an isolated vertex, so it consists of a path of length $d$ from $v$ and one
more step. If a vertex node $v$ were a cut node, a longest path from $v$
would lie in one part after deleting $v$, and a hyperedge at $v$ in another
part, together with one of its other vertices, would extend it to length
$d+1$.
\end{proof}

\begin{lemma}\label{lem:rigid}
Under the hypotheses of Lemma~\ref{lem:long}, every incidence-unicyclic
component is a loose $d$-cycle.
\end{lemma}

\begin{proof}
Let $Q$ be the cycle of the incidence graph. A hyperedge node off $Q$ would
lie in a tree hanging from $Q$, and the first vertex node on the way from
it to the rest of the component, which is either the attachment node or the
node following a hyperedge attachment node, would be a cut node separating
two parts that contain hyperedges. So every hyperedge lies on $Q$; its
third vertex is a pendant vertex, since a chord would create a second
cycle, and every vertex on $Q$ has degree two. The component is a loose
cycle with $L$ hyperedges, its longest Berge paths have length $L$, and
every vertex starts one; Lemma~\ref{lem:long} gives $L=d$.
\end{proof}

With two isolated vertices there is no linear tree component with a
hyperedge: one with at least two hyperedges has a vertex cut node, and a
single hyperedge has longest path length one. Hence
$|E(H)|\ge\lceil(n-2)/2\rceil$, and if $n$ is even and equality holds then
every non-isolated component is incidence-unicyclic, hence a loose
$d$-cycle, and $n=2td+2$.

\section{Proof of Theorem~\ref{thm:main}}

Let $H$ be Berge-$P_m$-saturated on $n\ge b$ vertices and $q=\lfloor n/b\rfloor$.
If a component is a single hyperedge, Lemma~\ref{lem:edge} applies. With no
isolated vertex, each linear tree component is itself Berge-$P_m$-saturated
on its vertex set and has at least four vertices, so by~(T) it has at least
$a$ hyperedges and at least $b$ vertices; there are at most $q$ of them, and
counting gives $2|E(H)|\ge n-q$. With one isolated vertex,
Corollary~\ref{cor:one} applies. With two isolated vertices
$|E(H)|\ge\lceil(n-2)/2\rceil$, which is at least $U(n)$ when $q\ge2$ and equals
$U(n)$ when $q=1$ and $n$ is odd. If $q=1$ and $n$ is even, so that
$b<n<2b$, then $U(n)=n/2$, and $|E(H)|=n/2-1$ forces disjoint loose
$d$-cycles and two isolated vertices, hence $n\equiv2\pmod{2d}$ and the
stated uniqueness. Conversely $\mathcal C_{d,t}$ attains $n/2-1$ at
$n=2td+2$ by Proposition~\ref{prop:cycles}, and~(C) supplies $U(n)$ hyperedges
at every other order. \qed

Exceptional orders exist for every $m\ge10$: one checks $a\ge d$, and
$t=\lceil a/d\rceil$ gives $a\le td\le 2a-1$, so $b<2td+2<2b$. They are the
orders $2td+2$ with $\lceil a/d\rceil\le t\le\lfloor (2a-1)/d\rfloor$.

\begin{center}
\begin{tabular}{rrrrll}
\toprule
$m$ & $d$ & $a$ & $b$ & exceptional $n$ & $\sat_3(n,\BP_m)$ versus $U(n)$\\
\midrule
10 & 8 & 15 & 31 & 34, 50 & 16, 24 versus 17, 25\\
11 & 9 & 19 & 39 & 56, 74 & 27, 36 versus 28, 37\\
12 & 10 & 23 & 47 & 62, 82 & 30, 40 versus 31, 41\\
13 & 11 & 27 & 55 & 68, 90 & 33, 44 versus 34, 45\\
\bottomrule
\end{tabular}
\end{center}

\section{Formal verification}\label{sec:lean}
\sloppy

The Lean~4.35.0 project in \texttt{lean/} uses Mathlib revision
\texttt{56a005855ea17e0b0958f73ec801ae140b0a1679}. A hypergraph is a
\texttt{Finset} of vertex sets; \texttt{IsBergePath} requires distinct
vertices, distinct hyperedges of the hypergraph and the linking condition
above; \texttt{ContainsBergePath H m} asks for one of length $m-1$; and
\texttt{IsBergePathSaturated} requires three-uniformity, freeness and a
Berge-$P_m$ after adding any absent three-element set. The saturation number
\texttt{satNumber n m} is the infimum of the sizes of saturated hypergraphs
on \texttt{Fin n}, and \texttt{Conjecture24ThreeUniform} states the
conjectured equality with $a$ and $U(n)$ written as above. The project
proves:
\begin{itemize}
\item \texttt{looseCycles\_saturated}: Proposition~\ref{prop:cycles} for
every $d\ge3$ and every $t$;
\item \texttt{witness34\_saturated} and \texttt{witness34\_card}: the
hypergraph on \texttt{Fin 34} with the triples
$\{16j+i,\,16j+((i+1)\bmod 8),\,16j+8+i\}$, $j\in\{0,1\}$, $0\le i\le7$, is
Berge-$P_{10}$-saturated and has $16$ hyperedges, while
\texttt{conjecturedSat 10 34 = 17};
\item \texttt{not\_conjecture24ThreeUniform},
\texttt{satNumber\_lt\_conjecturedSat} and
\texttt{conjecture24\_fails\_for\_every\_m}: the conjecture fails, the
saturation number is below $U(n)$ at every order $2t(m-2)+2$ strictly
between $b$ and $2b$, and every $m\ge10$ has such an order.
\end{itemize}
A module audit lists the transitive axioms of all 191 declarations of the
three proof modules, including generated ones; all depend only on
\texttt{propext}, \texttt{Classical.choice} and \texttt{Quot.sound}. The
lower bound of Theorem~\ref{thm:main}, the uniqueness statement and the
inherited results~(T) and~(C) are not formalized. A standard-library
Python program checks Proposition~\ref{prop:cycles} independently by
exhaustive search for $(m,n)=(10,34)$, $(10,50)$ and $(11,56)$.

\fussy
\section{Prior work and limitations}

The definitions, results~(T) and~(C), the lower bound of Theorem~2.3 and
Conjecture~2.4 are from the arXiv version of~\cite{EGKMS}; the journal
version~\cite{EGKMS-J} could not be accessed for this release, so whether
its wording of the conjecture differs has not been checked. Linearity of
Berge saturation numbers for small uniformities is due to English,
Gerbner, Methuku and Tait~\cite{EGMT}, and graph path saturation to
K\'aszonyi and Tuza~\cite{KT}. Searches in October 2026 of arXiv, journal
records and general web indexes did not find the loose-cycle construction,
a counterexample to Conjecture~2.4 or an exact formula; this does not
establish priority. Theorem~\ref{thm:main} covers only uniformity three,
$m\ge10$ and $n\ge b$.

\section*{Research process and AI disclosure}

This is fully autonomous agentic research. Dev D. Goyal is responsible
for publication. No expert validation or external peer review is
claimed, and the novelty assessment is provisional.

The question selection, the counterexample, the proof of
Theorem~\ref{thm:main}, the literature search and exhaustive computer
checks were produced by autonomous AI research agents in an earlier
session. Its records attribute the work, including an internal AI review
that rederived the formula and audited the proof, to an Astra agent
workflow, and they state that no Lean formalization was attempted.

Opus 5.5 by Anthropic used with Claude Code assessed the package for this
release. It read the relevant statements of~\cite{EGKMS}, checked the
written proof, rechecked the witnesses with its own exhaustive search,
wrote and ran the Lean formalization described in
Section~\ref{sec:lean}, and wrote this note. AI checks are not human peer
review.

\begin{thebibliography}{9}
\bibitem{EGKMS}
S.~English, N.~Graber, P.~Kirkpatrick, A.~Methuku and E.~C.~Sullivan,
\emph{Saturation of Berge hypergraphs},
\href{https://arxiv.org/abs/1710.03735v1}{arXiv:1710.03735v1} (2017).

\bibitem{EGKMS-J}
S.~English, P.~Gordon, N.~Graber, A.~Methuku and E.~C.~Sullivan,
\emph{Saturation of Berge hypergraphs}, Discrete Mathematics
\textbf{342} (2019), 1738--1761.
\href{https://doi.org/10.1016/j.disc.2019.01.031}{doi:10.1016/j.disc.2019.01.031}.

\bibitem{EGMT}
S.~English, D.~Gerbner, A.~Methuku and M.~Tait,
\emph{Linearity of saturation for Berge hypergraphs}, European Journal of
Combinatorics \textbf{78} (2019), 205--213.
\href{https://arxiv.org/abs/1807.06947}{arXiv:1807.06947}.

\bibitem{KT}
L.~K\'aszonyi and Zs.~Tuza, \emph{Saturated graphs with minimal number of
edges}, Journal of Graph Theory \textbf{10} (1986), 203--210.
\end{thebibliography}
\end{document}
